\section{Conclusion}
This work has demonstrated that both communication latency and computation latency exert a profound impact on the accuracy of cooperative perception, leading to severe safety risks in dynamic traffic environments. Beyond the widely acknowledged problem of communication latency, our study further proves that computation latency also plays a critical role in degrading perception accuracy. 
Through extensive evaluations on V2X-Sim and DAIR-V2X-C, we show that the proposed Dual-Latency Prediction Compensation Mechanism (DLPCM), together with the Vehicle-Specific Prediction Model (VSPM) and Model State Transmission (MST), achieves significant improvements in cooperative perception. The results demonstrate that our mechanism not only enhances real-time performance and prediction accuracy, but also alleviates communication pressure and reduces computational burden. The additional DAIR-V2X-C experiments further confirm the need for latency compensation in real cooperative perception and show that our VSPM and MST designs remain effective under real-scene BEV prediction, communication-state compression, and robustness evaluations. These findings validate the effectiveness and efficiency of the proposed framework in addressing the dual-latency challenge.
It is important to note that the emphasis of this work lies in the design of a new predictive compensation framework rather than the fine-tuning of specific prediction algorithms or compression strategies. While the framework already exhibits substantial advantages, future research will focus on improving the fidelity of prediction algorithms, exploring advanced compression methods for model state transmission, and conducting validations in large-scale heterogeneous cooperation. These directions will continue to strengthen the practicality of DLPCM and accelerate the safe and reliable deployment of cooperative perception in real-world intelligent driving systems.
